% \section{Nominal \titlemu HML}
In this section we present our logic \fhml.
We introduce the syntax and semantics of formulas, and show that the semantics is well defined. 
%For recursion, fixpoint operators are used and so we additionally show that least fixpoints exist. 
% We begin by providing the syntax for nominal \mhml\ formulas.
Recall the sets $\Atoms,\Sigma$ of {names} and tags, where the former is ranged over by $a, b,$ etc.
It will be useful to consider tags as having arbitrary finite {arities}, and not just unary as in FRAs, determined by a map $\ar:\Sigma\to\mathbb{N}$.
% ; the tags used in FRAs in particular have by definition arity 1.
Let us also assume a set $\var$ of \emph{value variables} and a set $\VAR$ of \emph{recursion variables}. These are ranged over respectively by
$x,y$, etc.\ and
$X,Y$, etc.
Recursion variables have also {arities}, given by a map $\ar:\VAR\to\mathbb{N}$ (note function overload). Here and in the sequel, when writing variable sequences $\vec x$ we shall tacitly assume they contain distinct elements.

\begin{definition}\label{def:mhmlr}
The formulas of  fresh HML with recursion (\fhml) are:
\begin{align*} 
  \text{Formulas} \ni\ \phi &::=\ u=u\mid \phi\lor\phi\mid \lnot\phi\mid \bigvee\nolimits_{x}\phi
  \mid \fresh{x}\phi\mid \diam{\ell}\phi\mid (\mu X(\vec x).\phi)(\vec u)\mid X(\vec u)\\
  \text{Values} \ni\ u &::=\ x \mid a\\
  \text{Labels} \ni\ \ell &::=\ (t,\vec u)\qquad (\text{with }\ar(t)=|\vec u|)
\end{align*}
where
in $X(\vec u)$ and $(\mu X(\vec x).\phi)(\vec u)$ we have $|\vec u|=|\vec x|=\ar(X)$.
\end{definition}


A variable $X$ can be free or bound, with binding performed with $\mu X$.
 As usually, we impose that each recursion variable $X$ appears within an even number of negation operations from its binder, implying that there must be at least one appearance of every such $X$.  
Value variables $x$ are bound by $\bigvee_{x}$, $\fresh{x}$ and $\mu X(\vec x)$. 
We say that a formula is \boldemph{firm} if it has no free value variables, and \boldemph{closed} if it has no free recursion variables.

The size of a formula is calculated in such a way that it be no smaller than the size of its support, and moreover, be linearly related to a concise machine representation of the formula (e.g.\ using de Bruijn indices~\cite{DEBRUIJN1972381}).


\begin{definition}\label{def:formSize}
    The size for each formula is inductively defined as:
  \begin{align*}
    |u=v| &= 2 &
    |\bigvee\nolimits_x\phi| = |\fresh{x}\phi| = |\lnot\phi| &= 1+|\phi|\\
    |\diam{t, \vec u}\phi| &= 1 + |\vec u| + |\phi| &
    | X(\vec u)| &= 1+|\vec u|\\
    |\phi\lor\psi| &= 1+|\phi| + |\psi| &
    | (\mu X(\vec x).\phi)(\vec u)|&= 1+|\phi|+2|\vec u|
    \end{align*}
\end{definition}

The semantics of formulas is given with respect to a given FLTS: a formula is mapped to a subset of its configurations. We shall use the following grammars for value and recursion variable substitutions respectively:
\begin{align*}
\gamma\ &::=\ \varepsilon\ \mid\ \gamma,\{\vec a/\vec x\} %\mid \theta[\sigma X(\vec x).\phi/X]
&
\theta\ &::=\ \varepsilon\  \mid\ \theta,\{\sigma X(\vec x).\phi/X\}
\end{align*}
Thus, substitutions are sequences of mappings which can be on sequences of distinct value variables or single recursion variables.
We may write $\dom(\gamma)$ for the union of the domains of all elements in $\gamma$, while in $\gamma,\{\vec a/\vec x\}$ we assume that $|\vec a|=|\vec x|$ and, for each $i$, $x_i\notin\dom(\gamma)$. A formula $\phi$ acted on by a substitution $\gamma$ is written $\phi\{\gamma\}$ and defined recursively by:
\begin{align*}
  \phi\{\varepsilon\} &= \phi &
  \phi\{\gamma,\{\vec a/\vec x\}\} &=  (\phi\{\gamma\})\{\vec a/\vec x\}
\end{align*}
Similar conventions and definitions relate to recursion variable substitutions.
% While value substitutions are maps, recursion variable substitutions are sequences of \emph{maplets}, i.e.\ singleton maps.
In addition, $\phi\{\theta\}$ is only ever going to be considered in cases where no variable capture be possible.

\begin{definition}\label{def:hdsemantics}
Given an FLTS $\mathcal{L}=\langle \mathcal{S},L,\rightarrow\rangle$ with $L\subseteq\{(t,\vec a)\in\Sigma\times\Atoms^*\mid \ar(t)=|\vec a|\}$, let us set\
$\mathcal{U}=\Pcal(\fstates)$\  and, for each $n \in \mathbb{N}$, \ 
    $\Ucal_n = \Atoms^n \to \Ucal$.\\
  A \emph{$\mathcal{U}$-variable assignment} is a finite partial map  $\xi:\VAR\rightharpoonup\bigcup_{n\in\mathbb{N}}\Ucal_n$ such that, for each $X\in\dom(\xi)$, $\xi(X)\in\mathcal{U}_{\ar(X)}$ ($=\Atoms^{\ar(X)} \to \Ucal$). Given a $\mathcal{U}$-variable assignment $\xi$, the semantics
$\sem{\phi}$
  of a firm formula $\phi$ with respect to $\xi$ is an element of $\Ucal$ given inductively by:
\begin{align*}
  \sem{a=b} &= \emptyset &
    \sem{\bigvee\nolimits_{x}\phi} &= \bigcup\nolimits_{a\in\Atoms}\sem{\phi\{a/x\}}\\
  \sem{a=a} &= \fstates &
    \sem{\fresh{x}\phi} &= \bigcup\nolimits_{a\#\phi,\xi} \{(s, H) \in \sem{\phi\{a/x\}} \mid a \notin H\}\\
  \sem{\phi_1\lor\phi_2} &= \sem{\phi_1}\cup\sem{\phi_2} &
  \sem{X(\vec a)} &= \xi(X)(\vec a)\\
  \sem{\lnot\phi} &= \fstates\setminus\sem{\phi} &
    \sem{(\mu X(\vec x).\phi)(\vec a)} &= (\lfp(\lambda f.\lambda\vec b. \sem[{\xi[X\mapsto f]}]{\phi\{\vec b/\vec x\}}))(\vec a)\\
      \sem{\diam{\ell}\phi} &\multicolumn{3}{l}{${}=\{(s, H)\in\fstates\mid \exists (s, H)\xrightarrow{\ell}(s', H').\,(s', H')\in\sem{\phi}\}$}
\end{align*}
When $\xi$ is empty, we may simply write $\sem[]{\phi}$.
\end{definition}

Note that the notation $\lambda\vec b. \sem{\phi\{\vec b/\vec x\}}$ is shorthand for the function $\{(\vec b,\sem{\phi\{\vec b/\vec x\}})\mid \vec b\in\Atoms^n\}$. 
More generally, the semantics is only defined on firm formulas.

For the semantics to be well-defined, we need to ensure that the required least fixpoints exist. From the Knaster-Tarski theorem,
this can be done by showing that $\Ucal_n$ is a complete lattice and that $\lambda f.\lambda \vec b.\sem[{\xi[X \mapsto f]}]{\phi\{\vec b / \vec x\}}$ is a monotone function. In particular, setting $F=\lambda f.\lambda \vec b.\sem[{\xi[X \mapsto f]}]{\phi\{\vec b / \vec x\}}$, we have\
$
\lfp(F) = \bigsqcap \{g\in \Ucal_n\mid 
F(g)\sqsubseteq g\}
$,
where order and least-upper-bounds in $\Ucal_n$ are given as expected:
\begin{align*}
f\sqsubseteq g &\iff \forall\vec a\in\Atoms^n.\,f(\vec a)\subseteq g(\vec a)
&
    \bigsqcap S &= \lambda\vec a.\,\bigcap\{f(\vec a)\mid f\in S \}
\end{align*}
for any $f,g\in\Ucal_n$ and $S\subseteq\Ucal_n$.

\begin{restatable}{lemma}{completelattice}\label{lem:completelattice}
% \begin{lemma}\label{lem:completelattice}
$\Ucal_n$ is a complete lattice and,
for any $\phi, \xi$,
 the function $\lambda f.\lambda \vec b.\sem[{\xi[X\mapsto f]}]{\phi\{\vec b/\vec x\}}$ is monotone.
% \end{lemma}
\end{restatable}

Finally, we show that for any \fhml\ formula $\phi$ and $\Ucal$-variable assignment $\xi$,  $\supp(\sem{\phi}) \subseteq \supp(\phi) \cup \supp(\xi)$. 
To do this, we show that the function $\lambda \phi, \xi.\sem{\phi}$ is equivariant, and then use the equivariance principle.

\begin{restatable}{lemma}{semanticsequivariant}\label{lem:semanticsequivariant}
    The function $\lambda \phi, \xi.\sem{\phi}$ is equivariant.
    Therefore, 
for any formula $\phi$ and $\Ucal$-variable assignment $\xi$, $\supp(\sem{\phi}) \subseteq \supp(\phi) \cup \supp(\xi)$.
\end{restatable}
% \begin{proof}
%     The proof is similar to Lemma~\ref{lem:equivariant}, except for the following additional rule. 
%     We show the new inductive step as follows:
%     \begin{itemize}
%         \item $\phi=\fresh{x}\phi'$:
%         In this case, $\semp{(\fresh{x}\phi')} = \bigcup_{a\in\Atoms}\{(s, H) \in \semp{(\phi'\{a/x\})} \mid a \notin H\}$. 
%         By inductive hypothesis, the latter is $\bigcup_{a\in\Atoms}\{(s, H) \in \pi \cdot \sem{(\phi'\{a/x\})} \mid a \notin H\}
%         % =\bigcup_{a\in\Atoms}\{\pi \cdot (s, H) \in \sem{(\phi'\{a/x\})} \mid \pi(a) \notin H\}
%         =\pi \cdot \bigcup_{a\in\Atoms}\{(s, H) \in \sem{(\phi'\{a/x\})} \mid a \notin H\}$.
%     \end{itemize}
% \end{proof}

% \begin{lemma}
%     For any formula $\phi$ and $\Ucal$-variable assignment $\xi$, $\supp(\sem{\phi}) \subseteq \supp(\phi) \cup \supp(\xi)$.
% \end{lemma}
% \begin{proof}
%     By the equivariance principle.
% \end{proof}

A direct consequence of the latter result is that $\sem{\phi}$ is finitely supported, and so is the function $\lambda f.\lambda \vec b.\sem[{\xi[X \mapsto f]}]{\phi\{\vec b/\vec x\}}$.

\begin{restatable}{theorem}{semanticswelldefined}\label{thm:semanticswelldefined}
The semantics of \cref{def:hdsemantics} is well defined. Its image is within the subset of $\Ucal$ of finitely-supported sets of configurations.
\end{restatable}

\begin{remark}[Derived connectives]\label{rem:derived}
  We can define standard dual connectives for \fhml\ as:
  \begin{align*}
    \phi\land\psi &\equiv \lnot(\lnot\phi\lor\lnot\psi) &
    \sq{\ell}\phi &\equiv \lnot\diam{\ell}\lnot\phi &                        (\nu X(\vec x).\phi)(\vec u) &\equiv \lnot(\mu X(\vec x).\lnot\phi\{\lnot X/X\})(\vec u)
  \end{align*}
  and write $u\neq v$ for $\lnot(u=v)$. The semantics of the above can be derived from \cref{def:hdsemantics} as expected. For instance:\
$
  \sem{(\nu X(\vec x).\phi)(\vec a)} = (\gfp(\lambda f.\lambda\vec b. \sem[{\xi[X\mapsto f]}]{\phi\{\vec b / \vec x\}}))(\vec a)
  $.\\
  In the same vein,
  we show below that fresh selection is its self-dual:
  $\sem{\lnot\fresh{x}\phi}=\sem{\fresh{x}\lnot\phi}$.
%   \begin{align*}
%     (s,H)\in\sem{\lnot\fresh{x}\phi} &
%     \iff (s,H)\notin\sem{\fresh{x}\phi} 
% \iff\forall a\#\phi,\xi.\,(s,H)\notin\sem{\phi\{a/x\}}\lor a\in H\\
% &\iff\forall a\#\phi,\xi,H.\,(s,H)\notin\sem{\phi\{a/x\}} \\
% &\implies\exists a\#\phi,\xi,H.\,(s,H)\notin\sem{\phi\{a/x\}}
% \iff\exists a\#\phi,\xi,H.\,(s,H)\in\sem{\lnot\phi\{a/x\}}\\
% &\implies\exists a\#\phi,\xi.\,(s,H)\in\{(s',H')\in\sem{\lnot\phi\{a/x\}}\mid a\not\in H'\}\\
% &\iff (s,H)\in\sem{\fresh{x}\lnot\phi}\\
% (s,H)\in\sem{\fresh{x}\lnot\phi} &
% \iff \exists a\#\phi,\xi,H.\,(s,H)\in\sem{\lnot\phi\{a/x\}}
% \iff \exists a\#\phi,\xi,H.\,(s,H)\notin\sem{\phi\{a/x\}}\\
% &\implies\exists a\#\phi,\xi,H.\forall b\#\phi,\xi,H.\,(s,H)=\swp{a}{b}\cdot(s,H)\notin\sem{\phi\{b/x\}}\\
% &\implies\forall b\#\phi,\xi,H.\,(s,H)\notin\sem{\phi\{b/x\}}
% \iff(s,H)\in\sem{\lnot\fresh{x}\phi}
%   \end{align*}
\end{remark}

As in~\cite{PittsGabbay,Pitts}, we can show that $\fresh{x}\phi$ can be taken as meaning ``for any fresh $x$, $\phi$ holds'' or ``there is some fresh $x$ such that $\phi$ holds''. 

\begin{lemma}
  For any $\phi,\xi$ and configuration $(s,H)$,
  \begin{align}
    (s,H)\in\sem{\fresh{x}\phi} 
    \iff \exists a\#\phi,\xi,H.\,(s,H)\in\sem{\phi\{a/x\}}\label{NEWone}\\
\iff \forall a\#\phi,\xi,H.\,(s,H)\in\sem{\phi\{a/x\}}\label{NEWtwo}
  \end{align}
Therefore, we have $\sem{\lnot\fresh{x}\phi}=\sem{\fresh{x}\lnot\phi}$.
\end{lemma}
\begin{proof}
  We note that \eqref{NEWone} is a restatement of the definition of the semantics for $\fresh{x}\phi$. For \eqref{NEWtwo}
it suffices to show the left-to-right inclusion, as the other one follows from there always being some $a\#\phi,\xi,H$.
  For any $(s,H)$:
  \begin{align*}
(s,H)\in\sem{\lnot\fresh{x}\phi} &
\iff\exists a\#\phi,\xi,H.\,(s,H)\in\sem{\phi\{a/x\}}\\
&\implies\exists a\#\phi,\xi,H.\forall b\#\phi,\xi,H.\,(s,H)=\swp{a}{b}\cdot(s,H)\in\sem{\phi\{b/x\}}\\
 &\implies \forall b\#\phi,\xi,H.\,(s,H)\in\sem{\phi\{b/x\}}
  \end{align*}
  For the final claim of this Lemma, for each $(s,H)$ we have:
    \begin{align*}
    (s,H)\in\sem{\lnot\fresh{x}\phi} &
    \iff (s,H)\notin\sem{\fresh{x}\phi} 
\iff\forall a\#\phi,\xi,H.\,(s,H)\notin\sem{\phi\{a/x\}}\\
&\iff\forall a\#\phi,\xi,H.\,(s,H)\in\sem{\lnot\phi\{a/x\}}
\iff (s,H)\in\sem{\fresh{x}\lnot\phi}
  \end{align*}
\end{proof}

\begin{remark}[Comparison with logics of Dam~\cite{Dam03} and Klin et al.~\cite{KlinL19,KlinHD}]
  The $\pi$-$\mu$-calculus is a recursive modal logic for $\pi$-calculus, accompanied with a proof system that is sound and complete for finitary processes~\cite{Dam03}.
% The focus in \emph{loc cit.} is less on decidability and more in proof systems.
  Its connectives are very similar to ours and, in fact, our syntax for the recursive connectives is taken from \emph{loc cit}. Being a logic specifically for the $\pi$-calculus, it employs connectives such as free/bound input and output, in addition to modal connectives that specify the communication channel. Bound I/O in particular makes fresh-name quantification redundant. 
  \\
  The $\mu$-calculi with atoms of~\cite{KlinL19,KlinHD} are the probably closest to our logic. While those logics only feature recursion variables of nullary arities, the vectorial version studied in~\cite{KlinL19} can capture scenarios where non-zero arities are needed. The history-dependent logic of~\cite{KlinHD} captures freshness with this connective and corresponding semantics (in our terminology):
  \begin{align*}
\phi &::= \cdots \mid \#v\quad \text{(for value $v$)}
    &
\sem{\#a} &= \{(s,H)\mid a\notin H\}
  \end{align*}      
The analogy is not fully exact as our semantics is over history-dependent LTSs, while in \emph{loc cit.} on nominal LTSs with the history being attached to the semantics function\,---\,but that should be sufficiently close for our discussion here. Using $\#u$, we can define fresh quantification:
  \[
\fresh{x}\phi \equiv \bigvee\nolimits_x\#x\land(x\neq v_1)\land\cdots\land(x\neq v_n)\land\phi
\]
where $v_1,\dots,v_n$ is an enumeration of all (free) values in $\phi$.
But one can be even more expressive. E.g.\ the following formula stipulates that there are exactly two names in the current history:
\[
  \bigvee\nolimits_{x,y}x\neq y\land\lnot(\#x\lor\#y)\land
  \bigwedge\nolimits_{z}z\neq x,y\implies\#z
\]
We find this added expressivity not necessary for specifying FLTSs built e.g.\ from FRAs, as the exact size of the history is in general not accessible. The weaker history-dependence coming from the fresh-quantifier allows for a somewhat simpler notion of bounded parity game in the next section, compared to~\cite{KlinHD}. 
\end{remark}

By reduction from satisfiability for the $\mu$-calculus with atoms~\cite{KlinL19}, we know that the satisfiability problem is undecidable for \fhml. In the next section we shall study instead the model checking problem, show it is decidable and calculate an upper bound for its complexity.
The following definition will be of use then.

\begin{definition}
	Given a \fhml\ formula $\phi$, its bounding depth $\|\phi\|$ is given recursively by:	
	\begin{align*}
		\|u = v\| &= \|X(\vec u)\| = 0  &\|\lnot \phi\| &= \|\diam{\ell}\phi\| = \|\phi\| &\|(\mu X(\vec x).\phi)(\vec u)\| = \|\phi\| + |\vec x|\\
		\|\phi \lor \psi\| &= \max(\|\phi\|, \|\psi\|) & \|\bigvee\nolimits_x \phi\| &= \|\fresh{x} \phi\| = 1 + \|\phi\|   
	\end{align*}
\end{definition}


% <<<<<<< HEAD
%\paragraph*{Examples revisited}
We conclude this section by revisiting \cref{example:freshPaths,example:sessions} and building formulas for them.
% where all paths (finite and infinite) have distinct names.

\begin{example} The property \textsc{\#All}
 can be represented by the formula $\lnot \phi$, where:
% =======
% \paragraph*{Examples revisited}
% We revisit Example~\ref{example:freshPaths}, where all paths (finite and infinite) have distinct names ($\#\textsc{SUT}$).
% This can be represented by the formula $\lnot \phi$, where:
% >>>>>>> refs/remotes/origin/main
\[
	\phi = \mu X.\bigvee\nolimits_x\diam{x}(X\lor\mu Y.\bigvee\nolimits_y \diam{y}(Y\lor x=y))
\]
Note we omit tags as they are not used here.
The subformula $\mu Y.\bigvee\nolimits_y \diam{y}(Y\lor x=y)$ represents all finite paths that eventually accept an $x$, that is, for each call of $Y$ we accept any name $y$ and either  we repeat $Y$ or, in case $x = y$, we stop. 
%
%In this case, $Y$ is a $\mu$-variable as we require the fixpoint stabilises once an action has been repeated.
Similarly, $\phi$ represents all finite paths where eventually, some $a$ is repeated, that is, for each call of $X$ we accept any name $x$ and either we 
repeat $X$, or $x$ is eventually accepted again (using $\mu Y.\bigvee\nolimits_y \diam{y}(Y\lor x=y)$).

Thus, $\lnot \phi$ represents all paths such that no name is ever repeated.
Note that, by negation elimination, $\lnot \phi$ can be written as:
$
\nu X.\bigwedge\nolimits_x\sq{x}( X\land \nu Y.\bigwedge\nolimits_y \sq{y}(Y\land x\neq y))
$.
\end{example}

\begin{example}
Let us consider \cref{example:sessions}, where we start, repeatedly use, and then terminate a session that cannot be resumed or reactivated. 
This can be represented by:
\[
	\psi = \nu X. \fresh{s}\diam{\mathsf{S}, s}.\mu Y.(\diam{\mathsf{U}, s}Y \lor \diam{\mathsf{T}, s}X)
\]
%\begin{align*}
%	\psi &= \fresh{s_0}\diam{\mathsf{S}, s_{0}}(\nu X(s_i).(\diam{\mathsf{U}, s_i}X(s_i) \lor \diam{\mathsf{T}, s_i}\fresh{s_{i+1}}\psi'))(s_0)
%	\\
%	\psi' &= \diam{\mathsf{S}, s_{i+1}}X(s_{i+1}) \land \lnot(\mu Y().\bigvee\nolimits_y \diam{y}(Y \lor s_{i} = y))
%\end{align*}
where $\mathsf{S,U,T}$ stand for start, use and terminate respectively.
% <<<<<<< HEAD
% We begin by selecting a fresh session $s_0$ that has not yet been used, and then starting it. 
% From there, we then enter the recursive variable $X$ with $s_0$ and then:
% \begin{itemize}
% 	\item use $s_0$, represented by accepting $s_0$ with an action tagged \emph{use}, and repeat $X$ with the same session,
% 	\item or terminate $s_0$ and select some new session $s_1$.  
% \end{itemize}
% In the latter case, we enter $\psi'$, where we start $s_1$ and repeat $X$ now with $s_1$, and we also ensure that the now terminated session $s_0$ is not used again in any capacity (as in the previous example).
% The use of the greatest fixpoint variable $X$ ensures that the property is never violated.
% \end{example}
% =======
We begin by entering the recursive $\nu$-variable $X$, where each time we enter $X$, we select a fresh session $s$ and perform the action $\mathsf{start}(s)$. 
From here, we enter the $\mu$-variable $Y$ where each time we enter, we repeatedly perform the action $\mathsf{use}(s)$ and eventually perform the action $\mathsf{terminate}(s)$. 
Once $s$ has been terminated, we repeat the process by re-entering the $\nu$-variable $X$, starting some new session $s'$. 
Note that the use of the $\mu$-variable $Y$ ensures that the session will eventually terminate, and the use of the $\nu$-variable $X$ ensures that the property is never violated.
\end{example}
%We begin by selecting a fresh session $s_0$ that has not yet been used, and then starting it. 
%From there, we then enter the recursive variable $X$ with $s_0$ and then:
%\begin{itemize}
%	\item use $s_0$, represented by accepting $s_0$ with an action tagged \emph{use}, and repeat $X$ with the same session,
%	\item or terminate $s_0$ and select some new session $s_1$.  
%\end{itemize}
%In the latter case, we enter $\psi'$, where we start $s_1$ and repeat $X$ now with $s_1$, and we also ensure that the now terminated session $s_0$ is not used again in any capacity (as in the previous example).
%The use of the greatest fixpoint variable $X$ ensures that the property is never violated.
%>>>>>>> refs/remotes/origin/main
%%% Local Variables:
%%% mode: latex
%%% TeX-master: "main"
%%% End:
